\chapter{Родительский гессиан калибровочно-замкнутого шумового модуля}
\label{ch:v8-gauge-closed-noise-parent-hessian-gate}

\section{Проверка корректности самого вопроса}

Предыдущий гейт построил минимальное полное шумовое факторпространство
\begin{equation}
 \mathcal N_{\rm full}
 =\mathcal N_{\rm tr}^{\rm gauge}\oplus\mathcal N_{\rm gauge},
 \qquad \dim_{\mathbb C}\mathcal N_{\rm full}=15+12=27.
 \label{eq:v8-full-noise-splitting-for-parent-test}
\end{equation}
Возникает естественный вопрос: наследует ли он уже проверенный гессиан
Тома VII с переходом $(7,0,20)\to(0,0,27)$?

Этот вопрос требует сначала сопоставить типы пространств. Гессиан Тома VII
задан на двадцатисемимерном \emph{вещественном срезе полей} $\delta A$.
Часть переноса нового шумового факторпространства имеет пятнадцать комплексных каналов,
то есть тридцать вещественных направлений после овеществление. Двенадцать
оставшихся каналов вообще являются операторами концевых секторов алгебры
наблюдаемых, а не вариациями конечной стрелки $A:\mathbb C^{11}\to
\mathbb C^{10}$.

\section{Два пространства переноса не совпадают}

Прямое сравнение методом SVD дало
\begin{align}
 \dim_{\mathbb R}T_{\rm VII}&=27,&
 \dim_{\mathbb R}\mathcal N_{\rm tr}^{\rm gauge,\mathbb R}&=30,\\
 \dim(T_{\rm VII}\cap\mathcal N_{\rm tr}^{\rm gauge,\mathbb R})&=23,&
 \dim(T_{\rm VII}+\mathcal N_{\rm tr}^{\rm gauge,\mathbb R})&=34.
 \label{eq:v8-field-noise-space-intersection}
\end{align}
Следовательно, старый срез содержит четыре направления, которых нет в новом
шумовом модуле переноса, а новый модуль содержит семь направлений, на которых
старый гессиан не был задан. Ни одно пространство не содержит другое.

Это не только счёт размерностей. При двенадцати конечных калибровочных пробах
максимальная норма выхода старого полевого среза равна
\begin{equation}
 1.47987812,
 \label{eq:v8-old-field-slice-gauge-leakage}
\end{equation}
то есть полевой срез, на котором доказана сигнатура Тома VII, сам не является
полным калибровочным представлением.

\section{Что всё же продолжается на новый модуль переноса}

Относительная Грам-часть действия определена для произвольной матрицы
$\delta A$ и потому допускает честное продолжение на тридцать вещественных
направлений переноса. Но её сигнатуры равны
\begin{equation}
 \operatorname{sig}H_{\rm Gram}^{0}=(30,0,0),
 \qquad
 \operatorname{sig}H_{\rm Gram}^{*}=(0,2,28).
 \label{eq:v8-relative-gram-on-full-transfer}
\end{equation}
Она не воспроизводит семимодовый запуск и оставляет две нулевые вакуумные
моды. Нужный селектор в Томе VII обеспечивала рёберно-ходжева часть, но она
определена на старом неинвариантном 27D-срезе и не имеет уже выведенного
продолжения на семь новых направлений.

Неоднозначность видна даже в скалярном контроле
\begin{equation}
 H_{\lambda}=\frac12H_{\rm Gram}+\lambda I_{30}.
 \label{eq:v8-undetermined-transfer-mass-completion}
\end{equation}
При $\lambda=0$ начало имеет тридцать отрицательных мод; при $\lambda=3$
получается $(6,14,10)$, при $\lambda=4$ --- $(2,0,28)$, а при
$\lambda\geq5$ отрицательных мод нет. Все $\lambda>0$ поднимают две
вакуумные нулевые моды. Значит, качественный запуск зависит от нового
правила продолжения и не наследуется автоматически.

\section{Почему квадратная матрица размера 27 была бы нетипизированной}

Двенадцать калибровочных шумов действуют как коммутаторы на
$M_{11}(\mathbb C)\oplus M_{10}(\mathbb C)$. Для их полевого гессиана нужны
пространственно-временные калибровочные потенциалы, кинетическая метрика, фиксация
калибровочных орбит и общий след с конечными полями переноса. Эти данные не получены
в Томе VII. Поэтому объединение пятнадцати комплексных каналов переноса и
двенадцати калибровочных каналов в одну численную матрицу гессиана смешало бы
координаты поля с операторами скачка.

\begin{theorem}[Несовпадение полевого и шумового носителей]
Калибровочно-замкнутое 27-мерное шумовое факторпространство не наследует единственный
родительский гессиан Тома VII. Овеществлённый модуль переноса и старый полевой
срез пересекаются в размерности 23, но имеют соответственно семь и четыре
собственных направления. Относительная Грам-часть продолжается, а
рёберный селектор и калибровочно-полевая часть --- нет.
\end{theorem}

\begin{nogocriterion}[Запрет формального 27D шумового гессиана]
Нельзя отождествлять число 27 шумовых каналов с 27 вещественными координатами
старого физического среза и переносить сигнатуру
$(7,0,20)\to(0,0,27)$ по совпадению размерностей. Сначала требуется единое
пространство полей суперсвязности, содержащее вариации переноса и калибровочные вариации,
после чего гессиан должен быть вычислен заново на факторпространстве по калибровочным орбитам.
\end{nogocriterion}

\begin{protocol}
\begin{enumerate}
 \item старый полевой срез: \textbf{27 вещественных направлений};
 \item полный шумовой модуль переноса: \textbf{15 комплексных измерений, 30 вещественных измерений};
 \item пересечение: \textbf{23 вещественных измерений};
 \item старые направления вне шумового модуля: \textbf{4};
 \item новые шумовые направления вне старого среза: \textbf{7};
 \item полная калибровочная инвариантность старого среза: \textbf{нет};
 \item Грам-гессиан на новом модуле: \textbf{вычислен};
 \item рёберно-ходжева продолжение: \textbf{не выведено};
 \item калибровочно-полевой гессиан: \textbf{не типизирован};
 \item старый локальный переход: \textbf{сохранён в прежней области};
 \item полный шумовой родитель гессиан: \textbf{не получен};
 \item следующий гейт: \textbf{единое калибровочно-замкнутое пространство
 полей суперсвязности}.
\end{enumerate}
\end{protocol}

Машинный сертификат сохранён в
\path{s2t_v8_gauge_closed_noise_parent_hessian_gate_results.json}.

\begin{thebibliography}{9}
\bibitem{v8-noise-parent-holevo}
A.~S.~Holevo,
\emph{Covariant Quantum Dynamical Semigroups: Unbounded Generators},
arXiv:quant-ph/9701037.
\bibitem{v8-noise-parent-gaussian-qms}
F.~Girotti, D.~Poletti,
\emph{Gaussian quantum Markov semigroups on finitely many modes admitting a
normal invariant state},
arXiv:2412.10020.
\end{thebibliography}